\begin{textsample}{POS Dim 1 – pi – Score 50.00 – pi/pi\_em\_13.txt}  \label{ex:f1_pos_082}
3.1.4 Metodologias de Ensino Etimologicamente, a palavra “ aluno ” provém do latim alumnus, \textbf{que} significa “ afilhado ”, “ discípulo ”, “ criança de peito ” e, este termo, é derivado de alere, \textbf{que} tem o significado de “ nutrir ” ou “ alimentar ”.
A ideia do termo, portanto, é de \textbf{que} o aluno é aquele \textbf{que} está sendo nutrido ou criado.
Nesse entendimento, aluno é aquele \textbf{que} se “ alimenta ” de conhecimento.
Com ênfase, este processo de “ alimentação ” precisa de outra pessoa, o professor.
O vínculo entre ambos se concretiza num espaço social de aprendizagem, a escola, em algumas áreas do conhecimento e no modelo da \textbf{pedagogia} ( ALUNO, 2020 ).
Daí \textbf{que} a metodologia tradicional de ensino, criada no \textbf{século} XVII, tem o professor ( aquele \textbf{que} “ nutre ” ou “ alimenta ” ) como centro do processo, detentor do conhecimento e responsável por transmiti -lo aos alunos. \textbf{Uma} metodologia pautada pela \textbf{disciplina} em \textbf{que} o conteúdo é o mais importante, predominando a aprendizagem pela \textbf{memorização}.
Porém, no mundo \textbf{atual}, diante das rápidas transformações sociais embaladas pelo avanço da ciência e das tecnologias, tem -se buscado superar este modelo tradicional pela compreensão de \textbf{que} não cabe ao professor “ alimentar ” seus alunos com o conhecimento, mas sim \textbf{despertar} neles a \textbf{vontade} de buscá -lo.
Pesquisas recentes, tanto nas áreas da educação quanto da \textbf{psicologia} e da neurociência, comprovam \textbf{que} “ o processo de aprendizagem é único e diferente para cada ser humano, e \textbf{que} cada \textbf{um} aprende o \textbf{que} é mais relevante e \textbf{que} faz sentido para si, o \textbf{que} gera conexões cognitivas e emocionais ” ( BACICH ; MORAN, 2018, p. 38 ).
Desse modo, o processo de aprendizagem varia para cada ser humano, segundo o ritmo de cada \textbf{um}, conforme suas habilidades, sua modalidade de aprendizagem e, também, pelas \textbf{influências} externas.
Isto permite a inferência de \textbf{que} o papel da escola perpassa a \textbf{simples} \textbf{tarefa} de ensinar conteúdos, pois, \textbf{enquanto} espaço social, ela também tem a missão de auxiliar os estudantes no seu autoconhecimento e no seu autodesenvolvimento, tornando -os responsáveis pelo ato de aprender, aprender a aprender, ensinar e ter autonomia intelectual.
Para tanto, a escola precisa se preparar para o cumprimento desta missão.
Conforme a BNCC, as aprendizagens essenciais devem \textbf{concorrer} para assegurar aos estudantes o desenvolvimento de dez competências gerais, \textbf{que} consubstanciam, no âmbito pedagógico, os direitos de aprendizagem e desenvolvimento.
Essas aprendizagens essenciais só se materializam mediante o conjunto de decisões \textbf{que} caracterizam o currículo em ação.
Para tanto, recomenda : > selecionar e aplicar metodologias e estratégias \textbf{didático-pedagógicas} diversificadas, recorrendo a ritmos diferenciados e a conteúdos complementares, se necessário, para trabalhar com as necessidades de diferentes grupos de alunos, suas famílias e cultura de origem, suas comunidades, seus grupos de socialização etc. ; ( BRASIL, 2018, p. 17 ).
Conhecimento ; Pensamento científico, crítico e criativo ; Repertório cultural ; Comunicação ; Cultura digital ; Trabalho e projeto de vida ; Argumentação ; Autoconhecimento e autocuidado ; Empatia e cooperação ; Responsabilidade e cidadania.
Estas são as dez competências gerais \textbf{que} todos os estudantes devem desenvolver durante a Educação Básica, de acordo com a BNCC ( BRASIL, 2018 ).
O grande desafio para a gestão do currículo nas redes e instituições de ensino é materializá -lo em ações e/ou estratégias adequadas \textbf{que} resultem em processos de aprendizagem capazes de contemplar o desenvolvimento dessas competências.
Para tanto, é \textbf{imprescindível} \textbf{que} toda e qualquer ação constante no planejamento curricular seja realizada de forma integrada e \textbf{que} haja \textbf{contextualização} e \textbf{interdisciplinaridade}.
Neste aspecto, Bacich e Moran ( 2018 ), ao discorrerem sobre as metodologias ativas de aprendizagem, definem -nas como estratégias de ensino \textbf{que} colocam o estudante no centro do processo de aprendizagem, por meio da utilização de diversos recursos ( principalmente tecnológicos ), na tentativa de experimentar diferentes formas de se aprender.
Elas compreendem a implantação de novas formas de ensino na prática escolar, modificando o modo como é ativado a aprender.
O objetivo da diversificação de práticas escolares e metodologias ativas de ensino é melhorar a gestão educacional, enriquecer o ensino, estimular a aprendizagem e abrir \textbf{horizontes} novos, tanto para alunos quanto para professores.
Para Beck ( 2018 ), a proposta de educação \textbf{que} ativa, estimula e envolve os estudantes é \textbf{bastante} antiga, mas o conceito sobre metodologias ativas é \textbf{que} surgiu recentemente.
Renomados \textbf{autores} como Freire, Dewey, Knowles, Rogers, Vygotsky não citam o termo “ metodologias ativas ”, porém, em suas obras, defendiam a aplicação dos princípios \textbf{que} as fundamentam.
Num passado remoto, a filosofia socrática ( \textbf{século} V a.C. ) já buscava ativar os ouvintes por meio de \textbf{um} \textbf{método} interrogativo.
Significa \textbf{dizer} \textbf{que}, se fôssemos buscar \textbf{um} ‘ idealizador ’, teríamos \textbf{que} voltar milênios na história da educação.
Todavia, é importante destacar \textbf{que}, “ na construção \textbf{metodológica} da Escola Nova, Dewey teve grande \textbf{influência} na ideia das metodologias ativas ao defender \textbf{que} a aprendizagem ocorre pela ação, colocando o estudante no centro dos processos de ensino e de aprendizagem ” ( DIESEL ; BALDEZS ; MARTINS, 2017, p. 272 ).
Segundo as análises de Paiva \textbf{et} al. ( 2016 ), as metodologias ativas promovem \textbf{uma} ruptura com o modelo tradicional de ensino e fundamentam -se em \textbf{uma} \textbf{pedagogia} problematizadora, em \textbf{que} o aluno é estimulado/desafiado a assumir \textbf{uma} postura ativa diante de sua jornada educativa, do seu aprendizado.
Num trabalho primoroso de buscar pontos de convergência entre as metodologias ativas de ensino e outras \textbf{abordagens} já consagradas do âmbito da (re)significação da prática docente, Diesel ; Baldezs ; Martins ( 2017 ) também explicam os princípios \textbf{que} constituem as metodologias ativas de ensino.
No entanto, não adentraremos no aprofundamento de cada \textbf{um} deles, por razões de não nos tornarmos repetitivos das ideias já desenvolvidas em \textbf{tópicos} anteriores ( e posteriores ) deste documento, mas é \textbf{oportuno} abordá -los para efeito de compreensão.
De toda maneira, são os seguintes os princípios \textbf{que} constituem as metodologias ativas de ensino : a ) Aluno : centro do processo de aprendizagem ; b ) Autonomia ; c ) \textbf{Problematização} da realidade e reflexão ; d ) Trabalho em equipe ; e ) Inovação ; f ) Professor : mediador, facilitador, ativador.
Embora com algumas congruências principiológicas quanto aos pressupostos \textbf{teóricos} e \textbf{metodológicos}, as metodologias ativas não são uniformes, pois há grande diversidade e diferentes modelos e estratégias para sua operacionalização, constituindo -se alternativas para o processo de \textbf{ensino-aprendizagem}, com diversos benefícios e desafios, nos diferentes níveis educacionais.
Dentre a gama de variedades, destacam -se : Aprendizagem \textbf{baseada} em projetos ; Aprendizagem \textbf{baseada} em problemas ; Gamificação ; Sala de aula invertida ; Aprendizagem entre pares ; Cultura Maker ; Storytelling ; Estudos do meio, etc. ( COLLOR, 2019 ).
Considerando a realidade \textbf{atual}, Bacich e Moran ( 2018 ) confirmam \textbf{que} dois conceitos são especialmente \textbf{poderosos} para a aprendizagem : aprendizagem ativa e aprendizagem híbrida.
As metodologias ativas têm foco no protagonismo do estudante, na sua autonomia frente ao aprendizado, no seu envolvimento direto, participativo e reflexivo em todas as etapas do processo, experimentando, criando, tendo o professor como mediador.
A aprendizagem híbrida tem foco na flexibilidade, na mistura e compartilhamento de espaços, tempos, atividades, materiais, técnicas e tecnologias \textbf{que} compõem esse processo ativo.
A junção de metodologias ativas com modelos flexíveis e híbridos tem apresentado significativas contribuições na solução para garantir a sociabilidade, encurtar a distância física entre professores e alunos.
No mundo \textbf{contemporâneo}, é notória a tendência já consolidada para modelos híbridos de aprendizagem escolar, nos quais os estudantes têm a combinação de atividades a distância ou remotamente por meio de vídeos e exercícios interativos, e encontros presenciais.
Mas vale destacar \textbf{que}, quanto às atividades a distância, o foco deverá estar menos centrado nas tecnologias e mais nas competências e habilidades de alunos e professores.
A Resolução nº 3/2018/MEC-CNE-CEB ( DCNEM ), em seu Art. 7º, § 2º, dispõe \textbf{que} : > O currículo deve contemplar tratamento \textbf{metodológico} \textbf{que} evidencie a \textbf{contextualização}, a diversificação e a \textbf{transdisciplinaridade} ou outras formas de interação e articulação entre diferentes campos de saberes específicos, contemplando vivências práticas e vinculando a educação escolar ao mundo do trabalho e à prática social e possibilitando o aproveitamento de estudos e o reconhecimento de saberes adquiridos nas experiências pessoais, sociais e do trabalho.
Conforme Bacich e Moran ( 2018 ), \textbf{um} currículo flexível oportuniza o protagonismo dos estudantes, estimula sua autonomia intelectual, seu autodidatismo, para \textbf{que} possam personalizar seu percurso conforme suas necessidades, possibilidades, expectativas, interesses e estilos de aprendizagem, conciliando a prática e a \textbf{teoria}.
Portanto, é dever das escolas adotar propostas \textbf{metodológicas} \textbf{que} elevem o aprendizado dos estudantes e garantam sua autonomia intelectual na interpretação do mundo pelas perspectivas real e virtual.
Mas é \textbf{preciso} compreender \textbf{que} transformar o aluno em protagonista do próprio aprendizado demanda, também, \textbf{uma} transformação na gestão educacional, haja vista \textbf{que} é necessário o planejamento \textbf{baseado} em evidências para se proceder com a implantação de estratégias \textbf{metodológicas} \textbf{que} façam com \textbf{que} o aluno não apenas receba o conhecimento entregue pelo professor, mas sim participe ativamente do processo.
Os estudantes do \textbf{século} XXI estão inseridos em \textbf{uma} sociedade do conhecimento, num mundo marcado pela aceleração e pela transitoriedade das informações de interconexão com a cultura digital.
Neste sentido, o processo de ensino e aprendizagem desenvolvido por meio de metodologias ativas apoiadas em tecnologias visa a superar as \textbf{abordagens} educacionais centradas na fala do professor e em atividades educativas \textbf{que} situam os estudantes em condições de expectadores passivos ( BACICH ; MORAN, 2018 ).
Não obstante, a incorporação das metodologias ativas e inovadoras na gestão escolar demanda revisão na formação dos professores \textbf{que}, por sua vez, também \textbf{necessitam} desenvolver suas próprias competências e habilidades para atuarem como protagonistas na autonomia de sua práxis, porque mais do \textbf{que} saber o \textbf{que} são e como funcionam, os docentes precisam se apropriar dessas metodologias.
Na visão de Moran ( 2015, p. 17 ), romper com os modelos tradicionais de ensino e trabalhar com as metodologias ativas \textbf{exige} grandes mudanças e tomadas de decisões quanto aos aspectos físicos, operacionais e organizacionais das instituições de ensino : > Todos os processos de organizar o currículo, as metodologias, os tempos, os espaços precisam ser revistos e isso é complexo, necessário e \textbf{um} pouco assustador, porque não temos modelos prévios bem sucedidos para aprender.
Estamos sendo pressionados para mudar sem muito tempo para testar.
Por isso, é importante \textbf{que} cada escola defina \textbf{um} plano estratégico de como fará estas mudanças. [ … ] Capacitar coordenadores, professores e alunos para trabalhar mais com metodologias ativas, com currículos mais flexíveis, com inversão de processos ( primeiro, atividades online e, depois, atividades em sala de aula ).
Podemos realizar mudanças incrementais, aos poucos ou, quando possível, mudanças mais profundas, disruptivas, \textbf{que} quebrem os modelos estabelecidos.
Ainda estamos avançando muito pouco em relação ao \textbf{que} precisamos.
Nesta acepção, compreende -se \textbf{que} a implementação do currículo piauiense nas redes e instituições de ensino prescinde de programas, projetos ou planos de formação continuada de professores, incluindo, dentre as temáticas, as metodologias de êxito e estratégias \textbf{didático-pedagógicas} diversificadas, adequadas às necessidades dos educandos em suas diferentes realidades.
Além disso, é necessário definir sobre o lugar \textbf{que} essas metodologias e estratégias \textbf{didático-pedagógicas} irão ocupar nas propostas pedagógicas das escolas para o efetivo alcance dos resultados em aprendizagens, considerando o desenvolvimento de competências e habilidades, o protagonismo e a autonomia dos estudantes no seu percurso formativo, articulado com o seu projeto de vida.
É importante destacar \textbf{que}, em se tratando de estratégias eficazes de ensino com o auxílio de recursos tecnológicos no processo de dinamização dos ambientes de aprendizagem e construção de novos saberes e, sobretudo, para atender ao desafio de levar educação de qualidade às mais longínquas comunidades, a Secretaria de Educação do Estado/SEDUC-PI conta com o Programa de Mediação Tecnológica Canal Educação, \textbf{que} tem como objetivo \textbf{qualificar} a oferta da educação básica, com mediação presencial, elevando o índice de escolarização, a inclusão social e o prosseguimento nos estudos dos estudantes piauienses.
Atualmente, o Programa garante \textbf{que} a oferta da educação básica nas diferentes modalidades, ministradas por professores habilitados, utilizando mídias educativas com transmissão via satélite, em regime presencial para 900 ambientes escolares instalados em áreas urbanas e rurais do Estado.
A interatividade é vivenciada pelo contato direto do \textbf{educando} do centro de recepção/ambientes escolares, com o professor ministrante no centro de Emissão ( estúdio ). \textbf{Detalhando} \textbf{um} pouco mais sobre a metodologia utilizada e operacionalização do programa, as aulas são transmitidas via satélite pela televisão e recepcionadas, em tempo real, pelas escolas com os kits tecnológicos, no formato de videoconferência.
As aulas são ministradas por professores especializados em cada área do conhecimento, em conjunto com professores mediadores e assistentes devidamente capacitados para orientar e aprofundar os conteúdos, a partir do estúdio instalado na TV Educativa do Estado, em Teresina-PI, resultando numa significativa interatividade entre professor e aluno.
Os alunos, organizados em turmas, assistem às aulas em recepção organizada, ou seja, estão presentes nos ambientes escolares de segunda à sexta-feira, nos horários estabelecidos para a transmissão das aulas.
Em cada sala de aula há \textbf{um} mediador pedagógico \textbf{que}, de forma também presencial, realiza a mediação das atividades pedagógicas, organiza o ambiente físico e educativo, \textbf{que} permite a realização das atividades/dinâmicas locais e a interação discente/docente ( ministrante ), por meio de \textbf{uma} webcam \textbf{que} transmite sua imagem, sua voz e dados, resultando num diálogo efetivo, ao vivo/tempo real, garantindo, assim, a completa comunicação/interatividade entre os \textbf{sujeitos} do processo de ensino e aprendizagem.
A metodologia vivenciada pelo programa exercita o imperioso respeito nos ambientes escolares destinado a ouvir e, ao apropriar -se da informação, fica o \textbf{educando} credenciado às críticas reflexivas fundamentando os debates posteriores e, também, à tomada de consciência do quanto se pode produzir intelectualmente ao aperfeiçoar -se o tempo destinado às atividades pedagógicas, \textbf{que} são os momentos destinados aos exercícios lançados pelo professor ministrante a serem trabalhados e discutidos pelo mediador pedagógico com os educandos e posterior retorno ao professor ministrante.
Relativamente ao alastramento da COVID-19 no mundo, caracterizado como estado de Pandemia, seguindo as deliberações da Organização Mundial de Saúde ( OMS ), bem como considerando as condições sanitárias do Estado do Piauí em relação à Pandemia, com observância ao protocolo 042 do Comitê de Operações Emergenciais ( COE/PI ) e após consulta à comunidade escolar sobre a impossibilidade da adoção do modelo híbrido, houve decretação da suspensão de aulas presenciais em toda a rede pública estadual.
Daí, por volta do mês de agosto de 2020 aos dias \textbf{atuais} ( abril/2021 ), as aulas passaram a ser ofertadas no formato exclusivamente remoto.
Diante desse contexto, com transmissão ao vivo pelo YouTube e TV aberta, o Programa de Mediação Tecnológica Canal Educação tornou -se a principal plataforma de conteúdo das aulas remotas da Rede Pública Estadual de Ensino do Piauí.
Na plataforma, estudantes e professores têm acesso aos materiais didáticos, além de aulas ao vivo, aulas gravadas e atividades complementares.
A estratégia foi pensada de forma \textbf{que} os estudantes da Rede continuem recebendo conteúdos dos diversos componentes curriculares em casa.
Considerando, pois, a implementação de alternativas pedagógicas com o uso de Tecnologias de Informação e Comunicação ( TIC’s ), a relevância do Programa é perceptível em \textbf{face} do impacto \textbf{imediato} com largo alcance social, com atendimento às situações emergenciais, à dinamização dos espaços e tempos de aprendizagem e à demanda reprimida, garantindo condições de acesso, permanência e elevação da escolaridade da população piauiense, combatendo, ainda, o elevado índice de evasão escolar.
Além disso, o Programa atende às necessidades da SEDUC-PI quanto à mobilização, formação continuada e capacitação técnica dos profissionais da educação, incidindo diretamente na elevação do padrão de qualidade educacional. \textbf{Ademais}, diante do avanço da tecnologia e das demandas da sociedade \textbf{contemporânea}, bem como das mudanças estruturais \textbf{preconizadas} pela Lei 13.415/2017, regulamentada pela Resolução nº 3/2018 ( DCNEM – MEC/CNE/CEB ), já se pode pensar numa ampliação da abrangência do Programa de Mediação Tecnológica Canal Educação, quanto à possibilidade de utilização desse instrumento na oferta da modalidade de EaD para o Novo Ensino Médio.
Em vista dos aspectos mencionados sobre concepções de metodologias de ensino, na perspectiva de garantir a formação humana e integral dos estudantes piauienses, a Rede Estadual de Ensino adota a concepção dos princípios de integração \textbf{metodológica} coerentes com o protagonismo, cuja \textbf{centralidade} de todo o processo de organização dessa aprendizagem ( estratégias \textbf{didático-pedagógicas} ) esteja, efetivamente, no estudante.
Com isto, são prevalentes os princípios das metodologias eficazes como possibilidades de ativar significativamente o aprendizado dos estudantes, colocando -os no centro do processo, com autonomia para a construção do próprio conhecimento articulado com o seu projeto de vida.

% matched lemmas: abordagem, ademais, atual, autor, basear, bastante, centralidade, concorrer, contemporâneo, contextualização, despertar, detalhar, didático-pedagógico, disciplina, dizer, educar, enquanto, ensino-aprendizagem, et, exigir, face, horizonte, imediato, imprescindível, influência, interdisciplinaridade, memorização, metodológico, método, necessitar, oportuno, pedagogia, poderoso, preciso, preconizar, problematização, psicologia, qualificar, que, simples, sujeito, século, tarefa, teoria, teórico, transdisciplinaridade, tópico, um, vontade
\end{textsample}
